% ============================================================
%  B题 多工序协同作业问题 —— 数学建模论文（前三问）
%  编译方式：xelatex B题_论文.tex （需要 ctex / xelatex 环境）
% ============================================================
\documentclass[12pt,a4paper]{article}

\usepackage[UTF8]{ctex}
\usepackage{geometry}
\geometry{left=2.5cm,right=2.5cm,top=2.8cm,bottom=2.8cm}

\usepackage{amsmath,amssymb,amsfonts}
\usepackage{booktabs}
\usepackage{multirow}
\usepackage{array}
\usepackage{longtable}
\usepackage{graphicx}
\usepackage{float}
\usepackage{caption}
\usepackage{subcaption}
\usepackage{enumitem}
\usepackage{listings}
\usepackage{xcolor}
\usepackage[hidelinks]{hyperref}
\usepackage{url}

% 表格内自动换行
\newcolumntype{P}[1]{>{\centering\arraybackslash}p{#1}}
\newcolumntype{L}[1]{>{\raggedright\arraybackslash}p{#1}}

\captionsetup{font=small,labelfont=bf}
\graphicspath{{./}}

% 代码样式
\lstset{
  basicstyle=\ttfamily\small,
  breaklines=true,
  frame=single,
  columns=fullflexible,
  showstringspaces=false,
  keywordstyle=\color{blue},
  commentstyle=\color{gray},
}

\title{\bfseries 多工序协同作业问题的建模与精确调度\\[0.4em]
\large ——基于关键路径与CP-SAT的全局最优调度模型}
\date{}

\begin{document}

\maketitle

\begin{abstract}
\noindent 
本文针对多工序协同作业问题，以最短总工期为目标，建立混合整数规划模型，综合运用关键路径法与 CP-SAT 精确算法求解四个逐步递进的子问题，并给出了可证明的全局最优调度方案。

\textbf{针对问题一}，本文利用 A 车间三道工序设备类型互不冲突的结构，将同型设备并联作业解释为等量均分工程量，把问题归结为以最小化完工时间为目标的确定性串行关键路径问题；通过解析计算各设备类型最短持续时间并推导完工时间下界 $C_{\max}\ge 200+1080+18000+18000=37280$ s，再由最早开始调度达到该下界，证明班组1独立完成 A 车间的最短时长为 \textbf{37280 s（10:21:20）}。

\textbf{针对问题二}，本文在问题一基础上将每台物理设备逐台编号，以工序起止时刻、同型设备投入台数、具体设备指派与设备访问弧为决策变量，建立带序列相关运输时间的混合整数规划模型，统一刻画工序链、同型并联、异型协同、设备容量与跨车间运输约束，并采用 OR-Tools CP-SAT 精确求解；CP-SAT 返回 OPTIMAL，可行解上界与分支定界严格下界均为 $144880$ s，故班组1完成 A--E 五车间的最短总工期为 \textbf{144880 s（40:14:40）}。

\textbf{针对问题三}，本文在问题二基础上引入班组2的设备，将每类设备可用数量翻倍并精确区分两班组不同初始位置，为 32 台物理设备分别建立以所属班组驻地为虚拟起点的访问回路，实现跨班组同型设备的协同与转移；CP-SAT 求得最优总工期 \textbf{73575 s（20:26:15）}，严格下界与之相等、相对间隙为 0，较问题二缩短 71305 s（降幅 49.22\%）。另编写不依赖 OR-Tools 的贪心—模拟退火程序，四个随机种子均独立收敛至同一最优值，构成交叉验证。

\textbf{针对问题四}，本文在问题三的双班组设备池上引入整数设备购置决策，建立“购置—调度”联合优化模型，采用“购置结构支配筛选＋设备负荷下界＋累积资源松弛下界＋CP-SAT 完整调度”的分解算法；在 500000 元预算约束下，最优方案为新增 1 台精密灌装机、3 台自动传感多功能机与 3 台高速抛光机，总费用恰为 500000 元，最短总工期 \textbf{32058 s（08:54:18）}，较问题三缩短 56.43\%；CP-SAT 上界与严格下界均为 32058 s，其余预算极大的竞争结构均被松弛下界排除，故为全局最优。

此外，本文对每一问均给出严格下界并证明全局最优性，配合独立启发式算法的交叉验证与脱离求解器的可行性复核，验证了模型与求解结果的可靠性。

\medskip
\noindent\textbf{关键词：}多工序协同作业；柔性作业车间调度；混合整数规划；CP-SAT 精确求解；关键路径法
\end{abstract}

\section{问题重述}

某工业制造系统需对 A、B、C、D、E 五个车间开展集中整修。每个车间的整修任务可抽象为由若干道工序组成的流程，各工序须按固定顺序依次执行；每道工序需占用特定类型的设备完成，部分工序还需两类设备协同作业。设备在同一车间内不同工序之间转移不计运输时间，跨车间转移则消耗由车间距离与设备移动速度共同决定的运输时间。

附件给出三张子表：

\begin{itemize}
  \item \textbf{工序流程表}：记录各车间工序的编号、所属车间、先后关系（按编号由小到大依次执行）及每道工序所需设备类型与对应作业效率（m$^3$/h）。
  \item \textbf{班组配置表}：记录班组1、班组2 的设备类型、设备编号、移动速度与设备单价。
  \item \textbf{车间距离表}：给出班组驻地与车间之间、车间与车间之间的距离。
\end{itemize}

综上，本题是一个资源受限的车间调度问题，其核心约束在于：工序顺序固定、部分工序需两类设备协同完成、设备跨车间复用需计入运输时间、同一时刻每台设备只能服务一道工序，且时间以整数秒计。

本文需依次解决以下四个问题：

\begin{enumerate}[label=\textbf{问题\,\arabic*：},leftmargin=2em]
  \item 仅由班组1独立承担 A 车间的全部整修任务，建立数学模型，求最短时长，并给出每台设备的工序安排。
  \item 仅使用班组1的设备完成 A、B、C、D、E 五个车间的整修任务，建立数学模型，求全部任务的最短时长及设备安排。
  \item 同时使用班组1、班组2 的设备完成五个车间任务，建立数学模型，求最短时长，并给出设备安排及所属班组。
  \item 在不超过 500000 元的设备购置预算下，综合确定设备购置方案与作业调度策略，使全部任务完成时间最短，并给出每台设备的调度安排及设备购买情况。
\end{enumerate}

\section{问题分析}

四个子问题由简到繁、逐层递进：问题一仅涉及 A 车间三道工序，问题二扩展到五车间并需考虑设备跨车间转移，问题三引入班组2 使可用设备翻倍，问题四进一步引入设备购置决策。本文据此逐步扩展建模口径，并对每一问给出可证明的全局最优解。

\subsection{问题一的分析}

问题一仅涉及 A 车间的三道工序 A1（缺陷填补）$\to$ A2（表面整平）$\to$ A3（强度检测），目标为最小化 A 车间的总完工时间。三道工序使用的设备类型互不重叠，因此对任一工序增加同型设备投入不会与其余工序争夺资源；于是问题可分解为“先求每道工序的最短持续时间，再按最早开始原则串接”，属于确定性串行关键路径问题，可用解析下界与关键路径法求得并证明全局最优解，无需启发式搜索\cite{pinedo}。

\subsection{问题二的分析}

问题二要求仅用班组1 的设备完成五车间全部任务。与问题一相比，同一台设备完成一个车间任务后可转移至另一车间继续工作，转移消耗运输时间。因此问题同时包含工序链约束、同型并联约束、异型协同约束、设备容量约束以及与访问顺序相关的跨车间运输约束，属于带序列相关转移时间的多资源受限作业调度问题（柔性作业车间调度\cite{jobshop}）。若仅按车间各自优化再取最大值会忽略共享设备在时空上的冲突；若只把同型设备视为总容量又无法刻画具体设备的转运路径。故须对每台物理设备逐台编号，将“设备投入数量—具体设备指派—访问顺序—起止时刻”纳入同一精确优化模型。

\subsection{问题三的分析}

问题三在问题二基础上引入班组2 的设备，每类设备可用数量翻倍，但两个班组的设备初始位置不同：每台设备首次执行任务须从其所属班组驻地出发，之后方可跨车间连续服务。因此问题三并非简单地把问题二工期减半，而是需同时决策设备投入规模、具体设备指派（含所属班组）、设备路径与工序先后关系的混合柔性作业车间调度问题。需对 32 台物理设备逐台编号并区分其初始距离，才能得到真正的最优协同方案。

\subsection{问题四的分析}

问题四要求在不超过 500000 元的设备购置预算下，同时确定各班组新增设备数量与五个车间的作业调度策略，使全部任务最终完工时间最短。与问题三相比，本问新增了整数购置决策；设备数量改变后，每道工序的并联台数、持续时间、具体设备分配与跨车间路线都会同步变化。因此该问题实质上是设备购置背包问题与带序列相关运输时间的柔性作业车间调度问题的耦合。若直接枚举全部“设备类型—班组”购置组合并逐一求调度，计算量较大；若仅依据问题三的设备利用率进行贪心采购，又可能忽略不同瓶颈设备之间的组合效应。故本文先通过资源负荷分析识别瓶颈，再采用“购置结构支配筛选＋累积资源松弛下界＋CP-SAT 完整调度”的分解策略，兼顾探索效率与结论可靠性。

\section{模型假设}

为保证模型的明确性与可计算性，本文统一作出如下假设：

\begin{enumerate}[label=\textbf{假设\,\arabic*：},leftmargin=1.6em]
  \item 各班组设备数量、单机作业效率、移动速度及距离数据均以附件为准，计算期间设备不发生故障或临时增减，作业效率保持不变。
  \item 同型设备性能完全一致；两个班组同一类型设备的作业效率相同。
  \item 同一工序投入多台同类型设备时，该类设备对应的总工程量在投入设备之间\textbf{等量均分}，各设备\textbf{同时开始}。
  \item 一道工序需要两类设备时，两类设备同时开始，各自完成本类设备对应的全部工程量；工序完成时刻取各类设备释放时刻的最大值，较早完成的设备可立即释放并投入后续任务。
  \item 每台设备在同一时刻至多服务一道工序，且一道工序一旦开始不可中断。
  \item 各车间内部工序严格按题目给定顺序执行；C 车间的 C3--C5 循环按三次展开为 C3-1、C4-1、C5-1、$\dots$、C5-3。
  \item 设备首次运输起点为其所属班组驻地；跨车间移动按附件距离计算，同一车间内转移时间忽略；设备完成最后一项任务后无须返回驻地。
  \item 不同设备可同时运输，运输过程不存在道路拥堵、装卸等待与相互干扰；装卸、调试等未给出的辅助时间忽略不计。
  \item 所有时间统一换算为整数秒；对除法产生的非整数秒\textbf{向上取整}，以保证离散秒制下的方案在连续时间意义下仍可行。
  \item 新购设备与同类型原设备的作业效率、移动速度与作业规则完全相同，购入后在整个计划期内均可使用。
  \item 购置预算为不超过 500000 元（不必全部支出）；在最短总工期相同的情况下，优先选择购置费用更低的方案。
\end{enumerate}

\section{符号说明}

为统一全文表述，定义如下符号（下标 $i$ 泛指展开后的工序，$k$ 泛指设备类型，$m$ 泛指具体设备）：

\begin{table}[H]
\centering
\caption{主要符号及其含义}
\small
\begin{tabular}{@{}llL{8.6cm}@{}}
\toprule
符号 & 类型 & 含义 \\
\midrule
$W=\{A,B,C,D,E\}$ & 集合 & 车间集合 \\
$I$ & 集合 & 展开后的全部工序集合（C3--C5 重复三遍） \\
$K$ & 集合 & 设备类型集合（自动化输送臂、工业清洗机、精密灌装机、自动传感多功能机、高速抛光机） \\
$M_k$ & 集合 & 第 $k$ 类设备的具体设备编号集合（问题二 $k$ 类共 $n_k$ 台；问题三为两班组同类设备之并） \\
$R_i$ & 集合 & 工序 $i$ 所需的设备类型集合 \\
$Q_i$ & 参数 & 工序 $i$ 的工程量（m$^3$） \\
$v_{ik}$ & 参数 & 第 $k$ 类单台设备执行工序 $i$ 的作业效率（m$^3$/h） \\
$n_k$ & 参数 & 班组1（或单班组）拥有的第 $k$ 类设备台数 \\
$d_{ab}$ & 参数 & 位置 $a$ 与位置 $b$ 之间的运输距离（m） \\
$s$ & 参数 & 设备移动速度，恒为 $2$ m/s \\
$\tau_{ab}$ & 参数 & 设备从位置 $a$ 转移到位置 $b$ 的运输时间（s） \\
$\tau^{0}_{g,w}$ & 参数 & 设备从班组 $g$ 驻地首次运至车间 $w$ 的时间（s） \\
$w(i)$ & 参数 & 工序 $i$ 所属车间 \\
$g(m)$ & 参数 & 设备 $m$ 所属班组 \\
$m_{ik}$ & 变量 & 工序 $i$ 投入的第 $k$ 类设备台数 \\
$S_i,\ C_i$ & 变量 & 工序 $i$ 的开始时刻、整体完成时刻 \\
$E_{ik}$ & 变量 & 工序 $i$ 中第 $k$ 类设备的释放（完成）时刻 \\
$x_{ikm}$ & 0-1 变量 & 工序 $i$ 的第 $k$ 类需求是否选用具体设备 $m$ \\
$y_{ikr}$ & 0-1 变量 & 工序 $i$ 对第 $k$ 类设备是否投入 $r$ 台 \\
$z_{ijm}$ & 0-1 变量 & 设备 $m$ 是否由任务 $i$ 直接转移至任务 $j$ \\
$c_r$ & 参数 & 第 $r$ 类设备的单价（元/台） \\
$z_{rg}$ & 整数变量 & 班组 $g$ 购买第 $r$ 类设备的台数 \\
$N_r^0,\ N_r$ & 参数 & 购置前、购置后第 $r$ 类设备的总台数 \\
$L_r$ & 参数 & 第 $r$ 类设备在全部工序上的理论负荷（台$\cdot$s） \\
$C_{\max}$ & 变量 & 全部任务的最晚完工时刻（总工期） \\
\bottomrule
\end{tabular}
\end{table}

\section{数据预处理}

\subsection{工序流程与设备需求}

附件“工序流程表”整理如下，其中“设备（效率 m$^3$/h）”列给出该工序所需设备类型及对应单机作业效率。C 车间的 C3、C4、C5 三道工序需按工艺顺序循环执行三次。

\begin{table}[H]
\centering
\caption{五个车间工序流程及设备需求}
\small
\renewcommand{\arraystretch}{1.05}
\begin{tabular}{llllll}
\toprule
车间 & 工序 & 工序名称 & 工程量/m$^3$ & 设备（效率 m$^3$/h） \\
\midrule
\multirow{3}{*}{A} & A1 & 缺陷填补 & 300 & 精密灌装机(200)、自动化输送臂(250) \\
& A2 & 表面整平 & 500 & 高速抛光机(100)、工业清洗机(250) \\
& A3 & 强度检测 & 500 & 自动传感多功能机(100) \\
\midrule
\multirow{4}{*}{B} & B1 & 除污清洗 & 120 & 工业清洗机(100) \\
& B2 & 灌浆填补 & 1500 & 精密灌装机(200)、自动化输送臂(300) \\
& B3 & 灌浆加固 & 360 & 精密灌装机(350) \\
& B4 & 抛光检测 & 360 & 高速抛光机(120)、自动传感多功能机(100) \\
\midrule
\multirow{5}{*}{C} & C1 & 清洗除锈 & 720 & 工业清洗机(250)、自动化输送臂(250) \\
& C2 & 精密灌浆 & 720 & 精密灌装机(350) \\
& C3 & 填补压实 & 360 & 精密灌装机(200)、自动化输送臂(250) \\
& C4 & 表面抛光 & 400 & 高速抛光机(120)、工业清洗机(100) \\
& C5 & 综合检测 & 400 & 自动传感多功能机(100) \\
\midrule
\multirow{6}{*}{D} & D1 & 清洗除锈 & 600 & 工业清洗机(250) \\
& D2 & 灌浆填补 & 800 & 精密灌装机(200)、自动化输送臂(300) \\
& D3 & 灌浆加固 & 450 & 精密灌装机(350) \\
& D4 & 抛光整平 & 1500 & 高速抛光机(120)、自动传感多功能机(300) \\
& D5 & 强度检测 & 1500 & 自动传感多功能机(300) \\
& D6 & 表面抛光 & 700 & 高速抛光机(100) \\
\midrule
\multirow{3}{*}{E} & E1 & 清洗除锈 & 1000 & 工业清洗机(250) \\
& E2 & 灌浆填补 & 600 & 精密灌装机(350) \\
& E3 & 检测修复 & 600 & 自动传感多功能机(300)、工业清洗机(100) \\
\bottomrule
\end{tabular}
\end{table}

注：C 车间 C3--C5 循环执行三次，展开后为 C3-1$\to$C4-1$\to$C5-1$\to$C3-2$\to$C4-2$\to$C5-2$\to$C3-3$\to$C4-3$\to$C5-3。各车间工序链为
\[
\begin{gathered}
A1\to A2\to A3,\qquad B1\to B2\to B3\to B4,\\
C1\to C2\to C3\text{-}1\to\cdots\to C5\text{-}3,\qquad
D1\to\cdots\to D6,\qquad E1\to E2\to E3 .
\end{gathered}
\]

\subsection{班组设备配置}

\begin{table}[H]
\centering
\caption{各班组的设备配置（移动速度均为 $2$ m/s）}
\small
\begin{tabular}{lccc}
\toprule
设备类型 & 班组1 台数 & 班组2 台数 & 设备编号示例 \\
\midrule
自动化输送臂 & 4 & 4 & 自动化输送臂$g$-j \\
工业清洗机 & 5 & 5 & 工业清洗机$g$-j \\
精密灌装机 & 5 & 5 & 精密灌装机$g$-j \\
自动传感多功能机 & 1 & 1 & 自动传感多功能机$g$-1 \\
高速抛光机 & 1 & 1 & 高速抛光机$g$-1 \\
\midrule
合计 & 16 & 16 & — \\
\bottomrule
\end{tabular}
\end{table}

\subsection{运输时间换算}

设设备移动速度为 $s=2$ m/s，则设备从位置 $a$ 到位置 $b$ 的运输时间向上取整为
\begin{equation}
\tau_{ab}=\left\lceil \frac{d_{ab}}{s}\right\rceil,\qquad s=2\ \text{m/s}.
\end{equation}
当设备连续执行同一车间的两项任务时 $a=b$，令 $\tau_{ab}=0$；首次出发使用班组驻地到目标车间的初始距离，之后使用车间间距离。据此可得运输时间表。

\begin{table}[H]
\centering
\caption{运输时间表（单位：s）}
\small
\begin{tabular}{@{}lccccc@{}}
\toprule
\multirow{2}{*}{位置} & \multicolumn{5}{c}{到各车间的初始运输时间} \\
\cmidrule(lr){2-6}
 & A & B & C & D & E \\
\midrule
班组1 驻地 & 200 & 310 & 230 & 355 & 200 \\
班组2 驻地 & 250 & 230 & 310 & 340 & 275 \\
\bottomrule
\end{tabular}

\medskip
\begin{tabular}{@{}lccccc@{}}
\toprule
车间 & A & B & C & D & E \\
\midrule
A & 0 & 510 & 525 & 450 & 700 \\
B & 510 & 0 & 550 & 815 & 360 \\
C & 525 & 550 & 0 & 260 & 425 \\
D & 450 & 815 & 260 & 0 & 515 \\
E & 700 & 360 & 425 & 515 & 0 \\
\bottomrule
\end{tabular}
\end{table}

\subsection{同型设备并联持续时间}

设工序 $i$ 的工程量为 $Q_i$，第 $k$ 类设备单机效率为 $v_{ik}$。若投入 $r$ 台同型设备并等量均分工程量，则每台设备承担 $Q_i/r$，该类型的持续工作时间为
\begin{equation}
p_{ikr}=\left\lceil\frac{3600\,Q_i}{r\,v_{ik}}\right\rceil,\qquad r=1,2,\dots,|M_k|.
\end{equation}
式 (2) 将“设备投入台数”与“工序持续时间”建立离散映射，使模型能够在“集中多机缩短单道工序”与“分散设备保障多车间并行”之间自动权衡。

\section{问题一：班组1 独立完成 A 车间的最优调度}

\subsection{问题一模型建立}

\subsubsection{初始运输时间}

班组1 驻地到 A 车间的距离为 $d_{1A}=400$ m，故初始运输时间为
\begin{equation}
\tau_{1A}=\left\lceil\frac{d_{1A}}{s}\right\rceil=\left\lceil\frac{400}{2}\right\rceil=200\ \text{s}.
\end{equation}
即全部设备于 00:00:00 同时出发，最早于 00:03:20 到达 A 车间，A1 的实际开始时间不得早于 200 s。

\subsubsection{同型多机并联模型}

对工序 $i$ 所需设备类型 $k$，若投入 $m_{ik}$ 台性能相同的设备，则各设备承担的工程量满足总量守恒
\begin{equation}
\sum_{j=1}^{m_{ik}} q_{ikj}=Q_i,\qquad 1\le m_{ik}\le n_k .
\end{equation}
第 $j$ 台设备完成其分配工程量所需的工作时间为
\begin{equation}
p_{ikj}=\left\lceil\frac{q_{ikj}}{v_{ik}}\times 3600\right\rceil .
\end{equation}
由于同型设备效率相同，为最小化该类型设备的最晚完成时间，应将工程量均匀分配，即 $q_{ikj}=Q_i/m_{ik}$，于是该类型完成工序 $i$ 的最短时间为
\begin{equation}
p_{ik}(m_{ik})=\left\lceil\frac{3600\,Q_i}{m_{ik}\,v_{ik}}\right\rceil .
\end{equation}
当工序同时需要多种设备类型时，每类设备均须完成 $Q_i$，故工序持续时间由最慢的设备类型决定：
\begin{equation}
p_i=\max_{k\in R_i}\bigl\{p_{ik}(m_{ik})\bigr\}.
\end{equation}

\subsubsection{工序先后关系与目标函数}

A 车间三道工序严格串行，其时间约束与目标函数为
\begin{equation}
S_{A1}\ge\tau_{1A},\qquad S_{A2}\ge C_{A1},\qquad S_{A3}\ge C_{A2},\qquad C_i=S_i+p_i ,
\end{equation}
\begin{equation}
\min\ C_{\max}=\min\ C_{A3}.
\end{equation}

\subsection{问题一模型求解}

由于 A1、A2、A3 使用的设备类型互不重叠，对任一工序增加同型设备投入不会与其他工序形成资源竞争，故在不设置设备启用成本的情况下，投入全部可用同型设备即可使各类型设备的完成时间最小；随后令每道工序在其前序工序完成后立即开始，即可得到最早开始调度。图 \ref{fig:p1flow} 给出求解流程。

\begin{figure}[H]
\centering
\includegraphics[width=0.72\textwidth]{e58e95ed08cc2d9934bdbe6ac4f6140f.png}
\caption{问题一建模与精确求解流程}
\label{fig:p1flow}
\end{figure}

\subsubsection{各工序持续时间}

按式 (6)--(7) 计算各设备类型作业时间：A1 同时投入 5 台精密灌装机与 4 台自动化输送臂；A2 投入 1 台高速抛光机与 5 台工业清洗机；A3 投入 1 台自动传感多功能机，结果见表 \ref{tab:p1calc}。

\begin{table}[H]
\centering
\caption{A 车间各工序设备配置及时间计算}
\small
\label{tab:p1calc}
\begin{tabular}{ccccccl}
\toprule
工序 & 工程量/m$^3$ & 设备类型 & 效率/(m$^3$/h) & 投入台数 & 单机分配/m$^3$ & 时间/s \\
\midrule
\multirow{2}{*}{A1} & \multirow{2}{*}{300} & 精密灌装机 & 200 & 5 & 60 & 1080 \\
 & & 自动化输送臂 & 250 & 4 & 75 & 1080 \\
\midrule
\multirow{2}{*}{A2} & \multirow{2}{*}{500} & 高速抛光机 & 100 & 1 & 500 & 18000 \\
 & & 工业清洗机 & 250 & 5 & 100 & 1440 \\
\midrule
A3 & 500 & 自动传感多功能机 & 100 & 1 & 500 & 18000 \\
\bottomrule
\end{tabular}
\end{table}

\subsubsection{最优性证明}

由表 \ref{tab:p1calc} 可见，A1 中两类设备均需 1080 s；A2 由唯一一台高速抛光机决定，持续 18000 s；A3 由唯一一台自动传感多功能机决定，持续 18000 s。现证明该方案全局最优：

首先，A1 投入班组1 全部 5 台精密灌装机与 4 台自动化输送臂时，两类设备的总处理能力分别为 $5\times200=1000$ m$^3$/h 与 $4\times250=1000$ m$^3$/h，故 A1 的持续时间不可能小于 $\lceil 300\times3600/1000\rceil=1080$ s。其次，班组1 仅有 1 台高速抛光机，A2 至少需要 $18000$ s；仅有 1 台自动传感多功能机，A3 至少需要 $18000$ s。又因 A1、A2、A3 严格串行且首次作业前至少需 200 s 运输时间，任意可行方案均满足
\begin{equation}
C_{\max}\ \ge\ 200+1080+18000+18000=37280\ \text{s}.
\end{equation}
本文构造的最早开始调度恰好在 37280 s 完成全部任务，达到式 (10) 给出的理论下界，故该方案是问题一在本文口径下的\textbf{全局最优解}。

\subsection{问题一结果分析}

全部设备于 00:00:00 从班组1 驻地出发，00:03:20 到达 A 车间，随后按 A1$\to$A2$\to$A3 连续执行，设备级调度方案见表 \ref{tab:p1schedule}（同型号设备完全同质，具体编号之间可互换，不影响最短总工期）。图 \ref{fig:p1gantt} 给出对应甘特图。

\begin{table}[H]
\centering
\caption{问题一最优设备调度方案}
\small
\label{tab:p1schedule}
\begin{tabular}{cclcccc}
\toprule
序号 & 设备编号 & 分配工程量/m$^3$ & 起始时间 & 结束时间 & 持续时间/s & 工序 \\
\midrule
1--5 & 精密灌装机1-1$\sim$1-5 & 60 & 00:03:20 & 00:21:20 & 1080 & A1 \\
6--9 & 自动化输送臂1-1$\sim$1-4 & 75 & 00:03:20 & 00:21:20 & 1080 & A1 \\
10 & 高速抛光机1-1 & 500 & 00:21:20 & 05:21:20 & 18000 & A2 \\
11--15 & 工业清洗机1-1$\sim$1-5 & 100 & 00:21:20 & 00:45:20 & 1440 & A2 \\
16 & 自动传感多功能机1-1 & 500 & 05:21:20 & 10:21:20 & 18000 & A3 \\
\bottomrule
\end{tabular}
\end{table}

\begin{figure}[H]
\centering
\includegraphics[width=0.95\textwidth]{fd0759c021467fc87c389290619e4eb5.png}
\caption{班组1 完成 A 车间的最优设备调度甘特图}
\label{fig:p1gantt}
\end{figure}

A1 于 00:03:20 开始、00:21:20 完成；A2 随即开始，由高速抛光机的 18000 s 作业决定，于 05:21:20 完成；A3 随后执行，于 10:21:20 完成。故班组1 完成 A 车间全部工序的最短总时长为
\begin{equation}
C_{\max}^{*(1)}=37280\ \text{s}=10:21:20 .
\end{equation}

\section{问题二：班组1 设备跨五车间精确调度}

\subsection{问题二模型建立}

本文对班组1 的 16 台物理设备逐台编号，建立以全厂最大完工时刻为主目标的混合整数规划模型\cite{mip}。设 $I$ 为展开后全部工序集合，$K$ 为设备类型集合，$M_k$ 为第 $k$ 类设备的具体编号集合，$R_i$ 为工序 $i$ 所需设备类型集合。

\subsubsection{设备投入数量与具体设备分配}

用 0-1 变量 $y_{ikr}$ 表示工序 $i$ 对第 $k$ 类设备是否选择投入 $r$ 台，用 $x_{ikm}$ 表示是否选用具体设备 $m$，二者通过设备数量等式相连：
\begin{align}
&\sum_{r} y_{ikr}=1, && \forall i\in I,\ k\in R_i,\\
&\sum_{m\in M_k} x_{ikm}=\sum_{r} r\,y_{ikr}, && \forall i\in I,\ k\in R_i,\\
&p_{ik}=\sum_{r} p_{ikr}\,y_{ikr}.
\end{align}
式 (12) 保证每类设备只选一个投入档位，式 (13) 将投入台数落实为具体设备编号，式 (14) 将离散设备数量选择映射为对应持续时间。

\subsubsection{工序同步与完成时刻}

同一工序所需的不同设备类型共享开始变量 $S_i$；第 $k$ 类设备完成自身工程量后的释放时刻为
\begin{equation}
E_{ik}=S_i+p_{ik},\qquad \forall i\in I,\ k\in R_i .
\end{equation}
工序的整体完成时刻取所有类型释放时刻的最大值：
\begin{equation}
C_i=\max_{k\in R_i}\{E_{ik}\}.
\end{equation}
这一写法能处理一类设备先完成、另一类继续工作的情况，避免把全部参与设备错误地占用到工序整体结束。

\subsubsection{车间内部工序约束}

设 $P$ 为附件规定的直接前后工序关系集合，则任一后继工序须在前序工序整体完成后方可开始：
\begin{equation}
S_j\ge C_i,\qquad \forall (i,j)\in P .
\end{equation}

\subsubsection{单台设备路径与运输约束}

对每台物理设备 $m$，建立以班组初始位置为虚拟仓库节点 0 的有向访问路径。若 $z_{ijm}=1$，表示设备 $m$ 完成任务 $i$ 后直接前往任务 $j$。每个被分配任务恰有一条进入弧与一条离开弧，未被分配的任务不进入该设备路径，流平衡关系为
\begin{equation}
\sum_j z_{ijm}=x_{im},\qquad \sum_j z_{jim}=x_{im}.
\end{equation}
当设备由任务 $i$ 转向任务 $j$ 时，任务 $j$ 的开始时刻必须覆盖设备在任务 $i$ 中的释放时刻与跨车间运输时间，采用指示约束写为
\begin{equation}
z_{ijm}=1\ \Rightarrow\ S_j\ \ge\ E_{im}+\tau_{w(i),w(j)} .
\end{equation}
若任务 $j$ 是设备 $m$ 从班组驻地出发后的第一项任务，则有
\begin{equation}
z_{0jm}=1\ \Rightarrow\ S_j\ \ge\ \tau_{0,w(j)} .
\end{equation}
模型通过 Circuit 全回路约束同时消除子环并决定每台设备的访问顺序。由此，同一设备不可能同时参与两道工序，且设备位置随其实际任务序列连续更新。

\subsubsection{目标函数}

令 $l(w)$ 表示车间 $w$ 的最后一道工序，则全厂总完工时间为
\begin{equation}
C_{\max}=\max_{w\in W}\{C_{l(w)}\}.
\end{equation}
第一阶段以最小总完工时间为主目标：
\begin{equation}
\min\ C_{\max}.
\end{equation}
为在多组等价最优排程中选择更紧凑的一组，第二阶段固定第一阶段得到的最优值 $C_{\max}^{*}$，再最小化全部工序完成时刻之和：
\begin{equation}
\min\ \sum_{i\in I} C_i,\qquad \text{s.t.}\ C_{\max}=C_{\max}^{*}.
\end{equation}
第二阶段不改变最短总工期，仅用于减少无意义的整体等待，使结果更便于实施与绘图展示。

\subsection{问题二模型求解}

上述模型同时包含 0-1 分配变量、离散设备数量变量、最大值关系与序列相关运输约束，属于组合调度问题。本文采用 OR-Tools CP-SAT 精确求解框架\cite{ortools}：将工序起止时刻设为整数变量，将设备分配与路径选择设为布尔变量，利用约束传播、冲突学习、线性松弛与分支定界搜索可行域。与仅返回近似解的启发式算法不同，求解器同时维护当前最好可行解与严格下界；当二者相等并返回 OPTIMAL 时即可判定全局最优。

为减少等价解，在班组内部对同型设备编号施加“低编号优先启用”的对称性消除约束。主目标达到严格最优后固定 $C_{\max}$，再执行紧凑化次级目标。求解流程如图 \ref{fig:p2flow} 所示。

\begin{figure}[H]
\centering
\includegraphics[width=0.72\textwidth]{图1_问题二精确建模与求解流程图.png}
\caption{问题二精确建模与求解流程}
\label{fig:p2flow}
\end{figure}

求解步骤如下：
\begin{enumerate}[label=\textbf{步骤\,\arabic*：},leftmargin=2em]
  \item 读取五车间工序链、工程量、设备效率、设备数量与距离矩阵，统一换算为整数秒。
  \item 建立工序层开始—完成变量、设备数量档位变量及具体设备分配变量。
  \item 对每台物理设备建立含虚拟仓库节点的访问回路，并加入与访问顺序相关的运输时间约束。
  \item 最小化 $C_{\max}$，记录当前最好可行值、严格下界与最优间隙。
  \item 固定最短总工期，最小化全部工序完成时刻之和，获得紧凑排程。
  \item 脱离求解器，独立复核工序链、设备不重叠、运输时间、并联工程量与最终完工时刻。
\end{enumerate}

\subsection{问题二结果分析}

精确模型返回状态 OPTIMAL，主目标最好可行值与严格下界均为 $144880$ s，相对最优间隙为 0，故班组1 完成 A--E 五个车间全部任务的最短总工期为
\begin{equation}
C_{\max}^{*(2)}=144880\ \text{s}=40:14:40 .
\end{equation}
各车间最终完成时刻见表 \ref{tab:p2finish}，其中 D 车间的 D6 在 144880 s 完成，决定全厂总工期。

\begin{table}[H]
\centering
\caption{问题二各车间最终完成时刻}
\small
\label{tab:p2finish}
\begin{tabular}{ccccc}
\toprule
车间 & 最后工序 & 完成时刻/s & 时分秒 \\
\midrule
A & A3 & 127540 & 35:25:40 \\
B & B4 & 109030 & 30:17:10 \\
C & C5-3 & 142465 & 39:34:25 \\
D & D6 & 144880 & 40:14:40 \\
E & E3 & 17237 & 04:47:17 \\
\bottomrule
\end{tabular}
\end{table}

表 \ref{tab:p2schedule} 给出一组达到最短总工期的紧凑工序级方案，“设备投入”列列出该工序实际调用的设备类型与台数；若同一工序需两类设备，两类设备在该工序开始时刻同时开始，工序结束时刻取两类设备完成时刻的较大者。图 \ref{fig:p2gantt} 给出班组1 的最优设备调度甘特图。

\begin{table}[H]
\centering
\caption{问题二最优工序调度方案}
\small
\label{tab:p2schedule}
\begin{tabular}{lllcccc}
\toprule
车间 & 工序 & 设备投入 & 开始 & 结束 & 持续/s \\
\midrule
A & A1 & 精密灌装机$\times$5、自动化输送臂$\times$4 & 01:42:05 & 02:00:05 & 1080 \\
A & A2 & 高速抛光机$\times$1、工业清洗机$\times$1 & 21:00:45 & 26:00:45 & 18000 \\
A & A3 & 自动传感多功能机$\times$1 & 30:25:40 & 35:25:40 & 18000 \\
\midrule
B & B1 & 工业清洗机$\times$1 & 00:05:10 & 01:17:10 & 4320 \\
B & B2 & 精密灌装机$\times$5、自动化输送臂$\times$4 & 03:57:56 & 05:27:56 & 5400 \\
B & B3 & 精密灌装机$\times$5 & 05:27:56 & 05:40:17 & 741 \\
B & B4 & 高速抛光机$\times$1、自动传感多功能机$\times$1 & 26:41:10 & 30:17:10 & 12960 \\
\midrule
C & C1 & 工业清洗机$\times$4、自动化输送臂$\times$4 & 00:03:50 & 00:47:02 & 2592 \\
C & C2 & 精密灌装机$\times$5 & 00:47:02 & 01:11:44 & 1482 \\
C & C3-1 & 精密灌装机$\times$5、自动化输送臂$\times$4 & 01:11:44 & 01:33:20 & 1296 \\
C & C4-1 & 高速抛光机$\times$1、工业清洗机$\times$1 & 01:33:20 & 05:33:20 & 14400 \\
C & C5-1 & 自动传感多功能机$\times$1 & 10:02:00 & 14:02:00 & 14400 \\
C & C3-2 & 精密灌装机$\times$5、自动化输送臂$\times$4 & 14:02:00 & 14:23:36 & 1296 \\
C & C4-2 & 高速抛光机$\times$1、工业清洗机$\times$2 & 17:32:00 & 20:52:00 & 12000 \\
C & C5-2 & 自动传感多功能机$\times$1 & 22:32:00 & 26:32:00 & 14400 \\
C & C3-3 & 精密灌装机$\times$5、自动化输送臂$\times$4 & 26:32:00 & 26:53:36 & 1296 \\
C & C4-3 & 高速抛光机$\times$1、工业清洗机$\times$5 & 29:50:20 & 33:10:20 & 12000 \\
C & C5-3 & 自动传感多功能机$\times$1 & 35:34:25 & 39:34:25 & 14400 \\
\midrule
D & D1 & 工业清洗机$\times$4 & 02:02:42 & 02:38:42 & 2160 \\
D & D2 & 精密灌装机$\times$5、自动化输送臂$\times$4 & 02:40:55 & 03:28:55 & 2880 \\
D & D3 & 精密灌装机$\times$5 & 03:28:55 & 03:44:21 & 926 \\
D & D4 & 高速抛光机$\times$1、自动传感多功能机$\times$1 & 04:57:40 & 17:27:40 & 45000 \\
D & D5 & 自动传感多功能机$\times$1 & 17:27:40 & 22:27:40 & 18000 \\
D & D6 & 高速抛光机$\times$1 & 33:14:40 & 40:14:40 & 25200 \\
\midrule
E & E1 & 工业清洗机$\times$4 & 00:54:07 & 01:54:07 & 3600 \\
E & E2 & 精密灌装机$\times$5 & 02:11:45 & 02:32:20 & 1235 \\
E & E3 & 自动传感多功能机$\times$1、工业清洗机$\times$3 & 02:47:17 & 04:47:17 & 7200 \\
\bottomrule
\end{tabular}
\end{table}

\begin{figure}[H]
\centering
\includegraphics[width=0.98\textwidth]{图2_班组1完成五个车间的最优设备调度甘特图.png}
\caption{班组1 完成五个车间的最优设备调度甘特图}
\label{fig:p2gantt}
\end{figure}

为验证结果的可靠性，本文另行编写不调用 OR-Tools 的贪心—模拟退火交叉验证程序\cite{sa}：以五条工序链的可行交织顺序为编码，采用串行调度生成机制枚举设备投入台数并选择具体设备，退火邻域为交换两个不同车间的优先位置。该程序独立收敛到同一最优值 $144880$ s，与精确模型结果一致，构成独立的可行性与数值交叉验证。

\section{问题三：双班组设备跨五车间精确调度}

\subsection{问题三模型建立}

问题三在问题二模型基础上引入班组2 设备。设 $g(m)\in\{1,2\}$ 为设备 $m$ 所属班组，$M_k$ 为两班组第 $k$ 类设备的并集（如高速抛光机共 2 台、自动传感多功能机共 2 台，其余三类各 10 台，合计 32 台）。决策变量、工序同步与车间内部约束与问题二一致，核心差异在于\textbf{首次运输约束}须按设备所属班组区分。

\subsubsection{设备投入规模与具体设备分配}

\begin{align}
&\sum_{r} y_{ikr}=1, && \forall i\in I,\ k\in R_i,\\
&\sum_{m\in M_k} x_{ikm}=\sum_{r} r\,y_{ikr}, && \forall i\in I,\ k\in R_i,\\
&p_{ik}=\sum_{r} p_{ikr}\,y_{ikr}.
\end{align}
班组编号已包含在具体设备 $m$ 中，因此同一工序既可只用一个班组，也可由两个班组的同型设备共同完成。

\subsubsection{工序同步与完成时刻}

\begin{equation}
E_{ik}=S_i+p_{ik},\qquad C_i=\max_{k\in R_i}\{E_{ik}\}.
\end{equation}

\subsubsection{车间内部工序先后关系}

\begin{equation}
S_j\ge C_i,\qquad \forall (i,j)\in P .
\end{equation}

\subsubsection{单台设备路径与运输约束}

对每台具体设备 $m$ 建立以其所属班组驻地为虚拟节点 0 的有向访问路径：
\begin{equation}
\sum_j z_{ijm}=x_{im},\qquad \sum_j z_{jim}=x_{im},
\end{equation}
\begin{equation}
z_{ijm}=1\ \Rightarrow\ S_j\ \ge\ E_{im}+\tau_{w(i),w(j)},
\end{equation}
\begin{equation}
z_{0jm}=1\ \Rightarrow\ S_j\ \ge\ \tau^{0}_{g(m),w(j)} .
\end{equation}
式 (32) 中的首次运输时间 $\tau^{0}_{g(m),w(j)}$ 取决于设备 $m$ 所属班组 $g(m)$，由此精确区分两个班组不同的初始距离。对称性消除约束仅在\textbf{同一班组}内部对同型设备编号施加，不同班组因初始位置不同不能跨班组对称化。

\subsubsection{目标函数}

\begin{equation}
C_{\max}=\max\{C_{A3},C_{B4},C_{C5\text{-}3},C_{D6},C_{E3}\},
\end{equation}
\begin{equation}
\min\ C_{\max}.
\end{equation}
主目标取得严格最优后，固定 $C_{\max}$，再最小化全部工序完成时刻之和以紧凑化排程，次级目标不改变最短总工期。

\subsection{问题三模型求解}

问题三同时包含 0-1 分配变量、离散设备数量变量、最大值约束和带序列相关运输时间的设备路径，本文继续采用 CP-SAT 精确求解。求解流程如图 \ref{fig:p3flow} 所示，步骤如下：

\begin{enumerate}[label=\textbf{步骤\,\arabic*：},leftmargin=2em]
  \item 读取工序、工程量、设备效率、两个班组设备编号与距离数据，统一换算为整数秒。
  \item 展开 C 车间重复工序，为每个工序创建开始、释放与完成变量。
  \item 枚举同型设备投入台数，建立数量选择变量与具体设备分配变量。
  \item 为每台设备建立含班组虚拟起点的访问路径，添加首次运输、跨车间运输与不重叠约束。
  \item 最小化 $C_{\max}$；达到严格最优后固定 $C_{\max}$，执行限时紧凑化次级目标。
  \item 脱离求解器，对导出方案逐项复核工序链、工程量、持续时间、设备冲突与运输时间。
\end{enumerate}

\begin{figure}[H]
\centering
\includegraphics[width=0.72\textwidth]{图1_问题三精确建模与求解流程图.png}
\caption{问题三精确建模与求解流程}
\label{fig:p3flow}
\end{figure}

\subsection{问题三结果分析}

\subsubsection{最短总工期}

CP-SAT 返回主目标状态 OPTIMAL，最优可行解与严格下界均为 $73575$ s，相对最优间隙为 0，故问题三最短总工期为
\begin{equation}
C_{\max}^{*(3)}=73575\ \text{s}=20:26:15 .
\end{equation}
与问题二仅使用班组1 时的 $144880$ s 相比，双班组协同缩短 $71305$ s，降幅约 $49.22\%$。各工序起止时刻见表 \ref{tab:p3schedule}，A、B、C、D、E 车间分别于 17:59:15、15:20:45、20:08:00、20:26:15、03:21:00 完成，D6 最晚结束，故 D 车间决定全厂总工期。

\begin{table}[H]
\centering
\caption{问题三最优方案的工序起止时刻}
\small
\label{tab:p3schedule}
\begin{tabular}{llcccc}
\toprule
车间 & 工序 & 开始时刻 & 结束时刻 & 工序持续时间/s \\
\midrule
A & A1 & 00:58:40 & 01:21:10 & 1350 \\
A & A2 & 10:42:20 & 13:12:20 & 9000 \\
A & A3 & 15:29:15 & 17:59:15 & 9000 \\
\midrule
B & B1 & 00:03:50 & 00:39:50 & 2160 \\
B & B2 & 02:26:32 & 03:11:32 & 2700 \\
B & B3 & 03:11:32 & 03:17:43 & 371 \\
B & B4 & 13:32:45 & 15:20:45 & 6480 \\
\midrule
C & C1 & 00:05:10 & 00:26:46 & 1296 \\
C & C2 & 00:26:46 & 00:39:07 & 741 \\
C & C3-1 & 00:39:07 & 00:49:55 & 648 \\
C & C4-1 & 00:49:55 & 02:49:55 & 7200 \\
C & C5-1 & 02:49:55 & 06:49:55 & 14400 \\
C & C3-2 & 06:49:55 & 07:00:43 & 648 \\
C & C4-2 & 08:53:35 & 10:33:35 & 6000 \\
C & C5-2 & 11:23:35 & 13:23:35 & 7200 \\
C & C3-3 & 13:23:35 & 13:34:23 & 648 \\
C & C4-3 & 15:11:55 & 16:51:55 & 6000 \\
C & C5-3 & 18:08:00 & 20:08:00 & 7200 \\
\midrule
D & D1 & 01:12:26 & 01:41:14 & 1728 \\
D & D2 & 01:41:14 & 02:05:14 & 1440 \\
D & D3 & 02:05:14 & 02:12:57 & 463 \\
D & D4 & 02:34:15 & 08:49:15 & 22500 \\
D & D5 & 08:49:15 & 11:19:15 & 9000 \\
D & D6 & 16:56:15 & 20:26:15 & 12600 \\
\midrule
E & E1 & 00:33:51 & 01:03:51 & 1800 \\
E & E2 & 01:03:51 & 01:21:00 & 1029 \\
E & E3 & 01:21:00 & 03:21:00 & 7200 \\
\bottomrule
\end{tabular}
\end{table}

\subsubsection{关键设备路径}

系统中高速抛光机每个班组仅 1 台，是最稀缺的关键资源之一。最优方案令两台抛光机协同执行 C4-1$\to$D4$\to$C4-2$\to$A2$\to$B4$\to$C4-3$\to$D6，并在相邻车间之间连续转运，见表 \ref{tab:p3critical}。

\begin{table}[H]
\centering
\caption{高速抛光机关键设备路径}
\small
\label{tab:p3critical}
\begin{tabular}{ccccc}
\toprule
工序 & 车间 & 开始 & 抛光作业结束 & 前序运输时间/s \\
\midrule
C4-1 & C & 00:49:55 & 02:29:55 & 初始至 C \\
D4 & D & 02:34:15 & 08:49:15 & C$\to$D：260 \\
C4-2 & C & 08:53:35 & 10:33:35 & D$\to$C：260 \\
A2 & A & 10:42:20 & 13:12:20 & C$\to$A：525 \\
B4 & B & 13:32:45 & 15:02:45 & A$\to$B：510 \\
C4-3 & C & 15:11:55 & 16:51:55 & B$\to$C：550 \\
D6 & D & 16:56:15 & 20:26:15 & C$\to$D：260 \\
\bottomrule
\end{tabular}
\end{table}

D6 于 16:56:15 开始，两台高速抛光机共同工作 12600 s，于 20:26:15 完成。该路径与自动传感多功能机承担的 C5、D4--D5、B4 和 A3 任务相互制约，构成决定总工期的关键资源链。图 \ref{fig:p3gantt} 给出双班组的最优设备调度甘特图。

\begin{figure}[H]
\centering
\includegraphics[width=0.98\textwidth]{图2_双班组完成五个车间的最优设备调度甘特图.png}
\caption{双班组完成五个车间的最优设备调度甘特图}
\label{fig:p3gantt}
\end{figure}

\subsection{算法交叉验证与最优解存在性}

为避免仅依赖单一求解器，本文另行编写不调用 OR-Tools 的贪心—模拟退火程序：以五条工序链的可行交织顺序为编码，采用串行调度生成机制枚举每道工序的设备投入台数并选择能最早到达相应车间的具体设备，退火邻域为交换两个不同车间的优先位置。贪心初解为 $86565$ s，使用四个不同随机种子，模拟退火均得到 $73575$ s，与精确模型完全一致，见表 \ref{tab:p3compare}。

\begin{table}[H]
\centering
\caption{问题三不同求解方法的结果比较}
\small
\label{tab:p3compare}
\begin{tabular}{lccc}
\toprule
方法 & 总工期/s & 求解状态 & 说明 \\
\midrule
CP-SAT 精确模型 & 73575 & OPTIMAL & 可证明全局最优 \\
贪心初解 & 86565 & 可行 & 用于退火初值 \\
模拟退火（种子 20260820/21/22/23） & 73575 & 启发式 & 均收敛至精确最优值 \\
\bottomrule
\end{tabular}
\end{table}

关于最优解的存在性与全局最优性：本题工序数量有限、每类设备数量有限，设备投入台数与设备访问顺序的组合均有限，且将全部任务依次执行即可构造至少一个可行解，故可行集合非空且离散有限，最小值必然能够取到。模拟退火四次独立收敛到 $73575$ s 构成独立可行性与数值交叉验证，但启发式算法不能独立证明不存在更短方案；严格最优性来自 CP-SAT 给出的
\begin{equation}
\text{可行解上界}\ 73575\ \text{s}=\text{分支定界严格下界}\ 73575\ \text{s}.
\end{equation}
故在当前建模口径下，$73575$ s 已被严格证明为问题三的全局最优总工期。

\section{问题四：预算约束下设备购置与精确调度}

\subsection{问题四模型建立}

问题四在问题三双班组设备池的基础上引入整数设备购置决策。企业最多投入 $500000$ 元购买新设备，新设备与同类型原设备具有相同的作业效率、移动速度与使用规则；需同时决定各类型设备的购买数量与所属班组、每道工序的投入台数、具体设备编号、设备访问顺序及工序起止时刻，使五个车间全部任务的最终完工时间最短。该问题不是单纯的采购背包问题，而是“整数购置决策＋多资源柔性作业车间调度＋序列相关运输”的联合组合优化问题。

\subsubsection{购置变量与预算约束}

设 $c_r$ 为第 $r$ 类设备单价，$z_{rg}$ 为班组 $g$ 购买第 $r$ 类设备的台数，$N_r^0$ 为问题三中两班组已有的第 $r$ 类设备总数。各类设备单价见表 \ref{tab:p4price}。

\begin{table}[H]
\centering
\caption{各类设备单价与原有台数}
\small
\label{tab:p4price}
\begin{tabular}{lccc}
\toprule
设备类型 & 单价 $c_r$/元 & 原有台数 $N_r^0$ \\
\midrule
自动化输送臂 & 50000 & 8 \\
工业清洗机 & 40000 & 10 \\
精密灌装机 & 35000 & 10 \\
自动传感多功能机 & 80000 & 2 \\
高速抛光机 & 75000 & 2 \\
\bottomrule
\end{tabular}
\end{table}

购置须满足预算约束
\begin{equation}
\sum_{r\in R}\sum_{g=1}^{2} c_r z_{rg}\le 500000,\qquad z_{rg}\in\mathbb{Z}_+ .
\label{eq:p4budget}
\end{equation}
购置后第 $r$ 类设备总量为
\begin{equation}
N_r=N_r^0+\sum_{g=1}^{2} z_{rg} .
\end{equation}
新购设备在购置时可自由指定所属班组，而班组归属只影响首项任务的出发距离，故可先把新设备视为“班组待定”设备，在其首项任务确定后归入到该首个车间初始运输时间更短的班组。这一变换不改变可行域，却能显著减少班组分配造成的对称解。

\subsubsection{设备台数选择与具体指派}

定义 $y_{irn}=1$ 表示工序 $i$ 对第 $r$ 类设备投入 $n$ 台，$x_{irm}=1$ 表示工序 $i$ 的第 $r$ 类需求选用具体设备 $m$。每个设备需求只能选择一个投入台数，且具体设备数量须与台数选择一致：
\begin{align}
&\sum_{n=1}^{N_r} y_{irn}=1, && \forall i\in I,\ r\in R_i,\\
&\sum_{m\in M_r} x_{irm}=\sum_{n=1}^{N_r} n\,y_{irn}, && \forall i\in I,\ r\in R_i,\\
&p_{ir}=\sum_{n=1}^{N_r} p_{ir}(n)\,y_{irn},
\end{align}
其中 $p_{ir}(n)=\left\lceil 3600 q_i/(v_{ir}n)\right\rceil$ 为投入 $n$ 台时的整数秒持续时间，已按式 (2) 预先计算，从而避免非线性除法直接进入求解器。

\subsubsection{工序先后与同步完成约束}

同一工序所需的不同设备类型共享开始变量 $S_i$，第 $r$ 类设备的释放时刻为 $E_{ir}=S_i+p_{ir}$，工序完成时刻取各类型释放时刻的最大值：
\begin{equation}
C_i=\max_{r\in R_i}\{E_{ir}\}.
\end{equation}
若 $j$ 为 $i$ 在同一车间的直接后继，则
\begin{equation}
S_j\ge C_i .
\end{equation}

\subsubsection{设备访问路径与运输约束}

对每台具体设备 $m$ 建立含虚拟仓库节点 0 的有向访问回路，令 $a_{mij}=1$ 表示设备 $m$ 完成任务节点 $i$ 后紧接着服务节点 $j$，$f_{mi}=1$ 表示 $i$ 为设备 $m$ 的第一项任务。回路流入、流出度与 $x_{irm}$ 联立，未选任务节点以自环跳过，已选节点恰有一条进入弧与一条离开弧。相邻任务间的时间约束为
\begin{equation}
a_{mij}=1\ \Rightarrow\ S_j\ge E_{ir}+\tau_{w(i),w(j)},
\end{equation}
首项任务的首次运输约束为
\begin{equation}
f_{mi}=1\ \Rightarrow\ S_i\ge \tau^{0}_{g(m),w(i)} .
\end{equation}

\subsubsection{目标函数（词典序优化）}

第一层主目标为最小化五个车间末道工序完成时刻的最大值：
\begin{equation}
\min\ C_{\max},\qquad C_{\max}=\max\{C_{A3},C_{B4},C_{C5\text{-}3},C_{D6},C_{E3}\}.
\end{equation}
主目标达到全局最优后固定 $C_{\max}$，再最小化购置费用
\begin{equation}
\min\ \sum_{r}\sum_{g} c_r z_{rg},
\end{equation}
最后固定前两层最优值，以 $\sum_i C_i$ 最小为排程紧凑化目标。由此形成“最短总工期—最低购置费用—非关键工序紧凑化”的词典序优化。

\subsection{问题四模型求解：分解算法与全局最优性证明}

直接对全部购置槽位与设备路径一次性交给求解器会产生大量等价班组分配。为此，本文采用“外层购置结构筛选—内层精确调度”的分解策略。

\subsubsection{设备负荷下界与瓶颈识别}

对设备类型 $r$，将其在全部工序上的理论设备工时累加，记为
\begin{equation}
L_r=\sum_{i\in I}\frac{3600\,q_i}{v_{ir}} .
\label{eq:p4loaddef}
\end{equation}
若购置后共有 $N_r$ 台该类设备，则任何调度均满足资源负荷下界
\begin{equation}
C_{\max}\ge \left\lceil\frac{L_r}{N_r}\right\rceil .
\label{eq:p4load}
\end{equation}
计算可知，高速抛光机的全部作业负荷为 $135000$ 台·秒，自动传感多功能机负荷为 $117360$ 台·秒，而问题三中两类设备均仅有 2 台，明显构成主导瓶颈。按单位费用带来的负荷下界改善幅度，初步购置结构集中于新增传感机与抛光机，并利用剩余预算补充精密灌装机。

\subsubsection{购置结构支配筛选与累积资源松弛下界}

若购置结构 $z'$ 在每类设备上的数量均不少于 $z$，则 $z'$ 的可行调度集合包含 $z$ 的全部调度，其最优总工期不可能更大，故任何未充分利用预算且仍能增加设备的结构均被某个预算极大结构支配。枚举预算内整数购置向量后，仅保留无法再购买最便宜设备的预算极大结构。

对第一轮筛选后仍可能优于当前上界的结构，建立累积资源松弛模型：保留车间工序链、可变投入台数与同型设备累计容量，但忽略具体设备编号、设备访问路径与跨车间运输。由于该模型删除了原问题的部分约束，其最优值必不大于原问题最优值，可作为严格下界：
\begin{equation}
\sum_{i\in A_r(t)} n_{ir}\le N_r,\qquad \forall r,t .
\label{eq:p4relax}
\end{equation}
其中 $A_r(t)$ 为时刻 $t$ 正在占用第 $r$ 类设备的工序集合。只要某购置结构的松弛下界不小于已知可行解上界 $32058$ s，该结构及其所有被支配结构均不可能产生更优方案。各预算极大竞争结构的松弛下界见表 \ref{tab:p4lb}。

\begin{table}[H]
\centering
\caption{预算极大竞争购置结构的累积资源松弛下界}
\small
\label{tab:p4lb}
\begin{tabular}{lccc}
\toprule
新增设备向量（臂/洗/灌/感/抛） & 费用/元 & 松弛下界/s & 判定 \\
\midrule
(0,0,1,3,3) & 500000 & 30208 & 进入完整模型，得 32058 s \\
(0,0,1,2,4) & 495000 & 32502 & 排除 \\
(0,1,0,2,4) & 500000 & 32480 & 排除 \\
(0,0,3,2,3) & 490000 & 33039 & 排除 \\
(0,1,2,2,3) & 495000 & 32732 & 排除 \\
(0,2,1,2,3) & 500000 & 32788 & 排除 \\
(1,0,1,2,3) & 470000 & 33127 & 排除 \\
(1,1,0,2,3) & 475000 & 32886 & 排除 \\
(2,0,0,2,3) & 485000 & 33230 & 排除 \\
\bottomrule
\end{tabular}
\end{table}

\subsubsection{完整 CP-SAT 求解与最优性证明}

购置向量按“自动化输送臂、工业清洗机、精密灌装机、自动传感多功能机、高速抛光机”排列。由表 \ref{tab:p4lb}，除 $(0,0,1,3,3)$ 外其余预算极大竞争结构的松弛下界均大于 $32058$ s，可直接排除。对唯一保留结构求解包含具体设备、完整运输与设备冲突的 CP-SAT 模型，求解器返回可行解上界 $32058$ s、严格下界 $32058$ s，状态为 OPTIMAL。为判断剩余 $35000$ 元是否必须用于精密灌装机，进一步对较低成本结构 $(0,0,0,3,3)$ 求解完整模型，其严格最优值为 $32125$ s，比 $32058$ s 慢 $67$ s，故达到最短工期所需的最低购置费用确为 $500000$ 元。图 \ref{fig:p4flow} 给出分解求解流程。

\begin{figure}[H]
\centering
\includegraphics[width=0.72\textwidth]{问题4_购置调度分解求解流程图.png}
\caption{问题四购置—调度分解求解流程}
\label{fig:p4flow}
\end{figure}

\subsection{问题四结果分析}

\subsubsection{设备购置方案}

最优方案新增 1 台精密灌装机、3 台自动传感多功能机与 3 台高速抛光机，总费用恰为 $500000$ 元，见表 \ref{tab:p4buy}。表 \ref{tab:p4buy} 给出与最终设备路线相匹配的一组班组归属；新设备在班组之间可能存在等价重新编号，但购置总量与最短总工期不变。

\begin{table}[H]
\centering
\caption{问题四最优设备购买方案}
\small
\label{tab:p4buy}
\begin{tabular}{lccccc}
\toprule
设备名称 & 班组1/台 & 班组2/台 & 单价/元 & 小计/元 \\
\midrule
自动化输送臂 & 0 & 0 & 50000 & 0 \\
工业清洗机 & 0 & 0 & 40000 & 0 \\
精密灌装机 & 1 & 0 & 35000 & 35000 \\
自动传感多功能机 & 0 & 3 & 80000 & 240000 \\
高速抛光机 & 3 & 0 & 75000 & 225000 \\
\midrule
合计 & 4 & 3 & — & 500000 \\
\bottomrule
\end{tabular}
\end{table}

\subsubsection{工序起止方案与最短总工期}

CP-SAT 返回 OPTIMAL，问题四最短总工期为
\begin{equation}
C_{\max}^{*(4)}=32058\ \text{s}=08:54:18 .
\end{equation}
与问题三不购置设备时的 $73575$ s 相比，缩短 $41517$ s，降幅约 $56.43\%$。各工序起止时刻见表 \ref{tab:p4schedule}，A、B、C、D、E 车间分别于 07:46:13、07:52:43、08:49:53、08:54:18、03:53:43 完成，D6 最晚结束，故 D 车间决定全厂总工期。

\begin{table}[H]
\centering
\caption{问题四最优方案的工序起止时刻}
\small
\label{tab:p4schedule}
\begin{tabular}{llcccc}
\toprule
车间 & 工序 & 开始时刻 & 结束时刻 & 工序持续时间/s \\
\midrule
A & A1 & 01:38:41 & 03:08:41 & 5400 \\
A & A2 & 05:07:51 & 06:22:51 & 4500 \\
A & A3 & 06:31:13 & 07:46:13 & 4500 \\
\midrule
B & B1 & 01:11:05 & 01:35:05 & 1440 \\
B & B2 & 02:33:42 & 03:33:42 & 3600 \\
B & B3 & 03:33:42 & 03:49:08 & 926 \\
B & B4 & 04:16:43 & 07:52:43 & 12960 \\
\midrule
C & C1 & 00:05:10 & 00:26:46 & 1296 \\
C & C2 & 00:26:46 & 00:38:00 & 674 \\
C & C3-1 & 00:38:00 & 00:48:48 & 648 \\
C & C4-1 & 00:48:48 & 01:36:48 & 2880 \\
C & C5-1 & 02:37:28 & 03:57:28 & 4800 \\
C & C3-2 & 03:58:18 & 04:09:06 & 648 \\
C & C4-2 & 04:09:06 & 04:59:06 & 3000 \\
C & C5-2 & 05:22:28 & 06:22:28 & 3600 \\
C & C3-3 & 06:22:28 & 06:35:58 & 810 \\
C & C4-3 & 06:35:58 & 07:35:58 & 3600 \\
C & C5-3 & 08:01:53 & 08:49:53 & 2880 \\
\midrule
D & D1 & 00:28:42 & 00:57:30 & 1728 \\
D & D2 & 00:57:30 & 01:24:10 & 1600 \\
D & D3 & 01:24:10 & 01:31:11 & 421 \\
D & D4 & 01:33:08 & 04:03:08 & 9000 \\
D & D5 & 04:03:08 & 05:18:08 & 4500 \\
D & D6 & 07:30:18 & 08:54:18 & 5040 \\
\midrule
E & E1 & 01:43:53 & 02:10:33 & 1600 \\
E & E2 & 02:10:33 & 02:27:42 & 1029 \\
E & E3 & 02:41:43 & 03:53:43 & 4320 \\
\bottomrule
\end{tabular}
\end{table}

\begin{figure}[H]
\centering
\includegraphics[width=0.98\textwidth]{问题4_最优设备调度甘特图.png}
\caption{问题四最优设备调度甘特图}
\label{fig:p4gantt}
\end{figure}

\subsubsection{探索算法与精确算法的对比}

为提供购置结构与工期上界，本文先用自适应大邻域搜索（ALNS）—模拟退火进行探索：以五条工序链的交织顺序为编码，采用离散事件串行调度生成器枚举投入台数并选择设备，邻域算子包括异车间交换、任务块搬移、拆除—修复及关键车间前拉。在“灌装机1、传感机3、抛光机3”结构下，ALNS 由 $43405$ s 的贪心初解改进至 $32355$ s，但仍比精确解慢 $297$ s（约 $0.93\%$），且启发式算法无法证明最优。本文因此转入精确分解算法，结果对比见表 \ref{tab:p4compare}。

\begin{table}[H]
\centering
\caption{问题四探索算法与精确算法结果对比}
\small
\label{tab:p4compare}
\begin{tabular}{lccc}
\toprule
方法 & 总工期/s & 相对差异 & 结论 \\
\midrule
ALNS—模拟退火 & 32355 & $+0.93\%$ & 近优可行解，无法证明最优 \\
购置结构分解＋CP-SAT & 32058 & 0 & 全局最优，上下界相等 \\
DAG 关键路径反算 & 32058 & 0 & 独立一致性验证通过 \\
\bottomrule
\end{tabular}
\end{table}

综上，问题四的全局最优性来自“可行解上界 $32058$ s $=$ CP-SAT 严格下界 $32058$ s”，而其余预算极大竞争结构均被累积资源松弛下界排除，故 $32058$ s 是全部预算可行购置与调度方案中的全局最短工期。

\section{模型的评价与改进}

\subsection{模型的优点}

\begin{enumerate}[label=\textbf{(\arabic*)}]
  \item \textbf{精确性与可证性}：问题二、三、四采用 CP-SAT 精确求解，同时维护可行解上界与严格下界，返回 OPTIMAL 且间隙为 0，能够在模型口径内\textbf{证明}全局最优，而非仅给出近似解。
  \item \textbf{建模完备性}：模型统一刻画了设备投入规模选择、具体设备指派、带序列相关运输时间的设备路径、异型设备协同、工序链约束与（问题四的）整数设备购置决策，完整还原了题目的全部关键约束。
  \item \textbf{口径一致与可复核}：全程采用整数秒、向上取整口径，并对导出的每一份方案脱离求解器独立复核工序链、工程量、持续时间、设备冲突与运输时间，保证连续时间意义下依然可行。
  \item \textbf{交叉验证}：以不依赖 OR-Tools 的贪心—模拟退火程序、关键路径法反算对精确解进行独立交叉验证，增强了结果的可信度。
  \item \textbf{分解与下界结合}：问题四通过“购置结构支配筛选＋设备负荷下界＋累积资源松弛下界”大幅压缩搜索空间，再对唯一保留结构做完整精确求解，兼顾了求解效率与全局最优性的严格证明。
\end{enumerate}

\subsection{模型的缺点与改进方向}

\begin{enumerate}[label=\textbf{(\arabic*)}]
  \item 本文按“同型设备等量均分工程量”的口径建立并联模型；若题目要求按其他分配方式（如允许非均分、允许设备中途退出），则需将 $x_{ikm}$ 与工程量分配变量联合建模，模型规模会增大。
  \item 问题二、三、四的 CP-SAT 求解依赖求解时间与内存；当工序规模进一步扩大时，可考虑分支定界与邻域搜索的混合策略，或在精确解基础上以元启发式求大规模近似。
  \item 模型未考虑设备故障、维护、人员换班及随机扰动等动态因素，后续可引入随机规划或鲁棒优化进行拓展。
\end{enumerate}

\section{结论}

本文针对多工序协同作业问题，分别建立了问题一的确定性关键路径模型与问题二、三、四的 CP-SAT 精确调度模型，得到并证明了四个子问题的最短总工期：\textbf{问题一} $37280$ s（10:21:20），\textbf{问题二} $144880$ s（40:14:40），\textbf{问题三} $73575$ s（20:26:15），\textbf{问题四} $32058$ s（08:54:18）。双班组协同相较单班组缩短约 $49.22\%$；在 500000 元预算下增购 1 台精密灌装机、3 台自动传感多功能机与 3 台高速抛光机后，又较问题三缩短约 $56.43\%$。全部结果均在整数秒、向上取整口径下通过独立可行性复核，并辅以模拟退火与关键路径法交叉验证，具有较强的可靠性与可解释性。

\begin{thebibliography}{9}
\bibitem{pinedo}
Michael L. Pinedo. \textit{Scheduling: Theory, Algorithms, and Systems}. Springer, 5th edition, 2016.
\bibitem{jobshop}
Brucker P. \textit{Scheduling Algorithms}. Springer, 5th edition, 2007.
\bibitem{mip}
Wolsey L A, Nemhauser G L. \textit{Integer and Combinatorial Optimization}. Wiley-Interscience, 1999.
\bibitem{ortools}
Laurent Perron, Vincent Furnon. OR-Tools CP-SAT Solver. Google. \url{https://developers.google.com/optimization/cp/cp_solver}, 2024.
\bibitem{sa}
Kirkpatrick S, Gelatt C D, Vecchi M P. Optimization by simulated annealing. \textit{Science}, 1983, 220(4598): 671--680.
\end{thebibliography}

\appendix
\section{求解程序说明}

本文求解程序均基于 Python 编写，关键程序文件与输出对应关系如下：

\begin{itemize}
  \item \texttt{问题1\_最优调度.py}：问题一关键路径解析求解，输出 37280 s。
  \item \texttt{问题2\_最优调度.py}：问题二 CP-SAT 精确求解，输出 144880 s；辅以 \texttt{问题2\_模拟退火交叉验证.py}。
  \item \texttt{问题3\_最优调度.py}：问题三双班组 CP-SAT 精确求解，输出 73575 s；辅以 \texttt{问题3\_模拟退火交叉验证.py}。
  \item \texttt{问题4\_预算约束下最优购置与调度.py}：问题四购置—调度联合 CP-SAT 精确求解，输出 32058 s；\texttt{问题4\_全局最优分解求解.py}：购置结构分解与全局最优性证明；\texttt{问题4\_关键路径法独立交叉验证.py}：独立交叉验证。
  \item 结果表分别对应 \texttt{问题2\_最优设备调度方案.csv}、\texttt{问题3\_最优设备调度方案.csv}、\texttt{表3\_问题3结果.csv}、\texttt{表4\_问题4结果.csv} 及 \texttt{表5\_问题4设备购买情况.csv}。
\end{itemize}

以问题一为例，其核心求解逻辑为：对每道工序的每种所需设备类型，投入全部可用同型设备并等量均分工程量，取同一工序内各设备类型持续时间的最大值为工序持续时间，再沿 A1$\to$A2$\to$A3 串接最早开始时刻，即可得到并证明全局最优解。

\begin{lstlisting}[language=Python,caption={问题一核心求解逻辑（节选）},label={lst:p1}]
def ceil_seconds(workload_m3, efficiency_m3_per_hour):
    return math.ceil(workload_m3 / efficiency_m3_per_hour * 3600)

# 对每道工序的每种设备类型：投入全部可用同型设备，均分工程量
for operation in OPERATIONS:
    selected = []
    for requirement in operation.requirements:
        chosen = list(equipment[requirement.equipment_type])  # 全部投入
        workload_per_machine = operation.workload_m3 / len(chosen)
        duration = ceil_seconds(workload_per_machine,
                                requirement.efficiency_m3_per_hour)
        for machine in chosen:
            selected.append((requirement, machine, workload_per_machine, duration))
    # 同一工序所需设备同时开始，并等待前序工序完成
    operation_start = max(predecessor_finish,
                          max(m.available_time for _, m, _, _ in selected))
    ...
    predecessor_finish = operation_finish  # 串接下一道工序
\end{lstlisting}

\end{document}
